\newpage
\section{Experiments}
\setcounter{subsection}{-1}

% \subsection{Overview and scientific goals}
% In this chapter I present a hypothesis-driven experimental program studying transfer learning in time-series forecasting using TimesFM as the primary foundation model. I prioritize reproducible evaluation on public benchmarks and compute-aware adaptation methods. I do not attempt to recreate proprietary pretraining assets (e.g., Google Trends), but instead quantify transfer behavior using open checkpoints and public datasets \cite{timesfm_das2024,timesfm_repo,timesfm_hf_v25}.

% \subsection{Assumptions and constraints}
% I assume a solo research setting with limited or uncertain compute. I therefore define three compute tiers: (i) low: single GPU (e.g., RTX 3060 / T4) or CPU-only, (ii) medium: 2--4 GPUs, (iii) high: cluster access (e.g., A100). I report both performance and compute (GPU-hours, peak memory), using the official TimesFM system requirements as reference for inference-scale memory/speed \cite{timesfm_sysreq}.
% I assume I do not have access to proprietary Google Trends pretraining data; I therefore rely on open checkpoints and non-leaking public datasets for any continued-pretraining experiments \cite{timesfm_hf_v25,gifteval_aksu2024}.

% \subsection{Datasets}
% I evaluate on the following public benchmarks:
% (i) the Monash forecasting repository \cite{monash_godahewa2021},
% (ii) Darts benchmark datasets and tooling \cite{darts_herzen2022},
% (iii) the ETT benchmarks from ETDataset \cite{etdataset_zhou},
% (iv) GIFT-Eval for modern foundation-model evaluation, and GiftEvalPretrain for leakage-free continued pretraining \cite{gifteval_aksu2024}.

% \subsection{Models, metrics, and statistical tests}
% \paragraph{Models.}
% My primary models are TimesFM v1.0 and v2.5 open checkpoints \cite{timesfm_hf_v1,timesfm_hf_v25}. For context, I include Chronos as an external pretrained probabilistic baseline \cite{chronos_ansari2024}.

% \paragraph{Metrics.}
% For point forecasting I report MAE and MASE \cite{hyndman2006mase}. For probabilistic forecasting I report CRPS and quantile loss (pinball/WQL) and assess calibration via empirical coverage and sharpness \cite{gneiting2007scoring}.

% \paragraph{Statistical testing.}
% For uncertainty estimates I use paired bootstrap confidence intervals. For paired forecast comparison I use the Diebold--Mariano test on matched loss differentials, with multiple-comparisons correction where applicable \cite{diebold1995dm}.

% \subsection{Experiment suite (E1--E10)}
% \paragraph{E1 (Zero-shot replication).}
% I replicate a paper-aligned zero-shot evaluation across Monash/Darts/ETT. For Monash I follow the TimesFM paper and filter datasets with missing values (resulting in 18 datasets) and report scaled MAE relative to a naive last-value baseline \cite{timesfm_das2024}. For ETT I report horizons 96 and 192, using both rolling windows (stride 1) when feasible and a compute-controlled disjoint-window evaluation (stride set to horizon) as in the TimesFM long-horizon benchmark documentation \cite{timesfm_longhorizon}.

% \paragraph{E2 (Context-length sweep).}
% I quantify sensitivity to context length (forecasting-as-prompting) by sweeping context lengths from 64 up to the model’s supported long-context limit, holding other settings fixed \cite{timesfm_repo,timesfm_hf_v25}.

% \paragraph{E3 (Robustness benchmark).}
% I apply controlled corruptions to the context only (missingness, block missingness, outliers, level shifts) and measure degradation relative to clean performance.

% \paragraph{E4 (Probabilistic calibration).}
% Using TimesFM 2.5 quantiles, I evaluate raw quantiles versus conformalized quantile regression (CQR) calibration \cite{romano2019cqr}, isotonic calibration, and head-only fine-tuning of uncertainty heads. I report CRPS and calibration diagnostics \cite{gneiting2007scoring}.

% \paragraph{E5 (Covariates via XReg).}
% I evaluate covariate integration using TimesFM’s documented XReg approach, comparing ``TimesFM + XReg'' and ``XReg + TimesFM'' under explicit covariate-availability regimes to prevent leakage \cite{timesfm_xreg}.

% \paragraph{E6 (Residual adapters).}
% I train lightweight residual correction models (ridge/GBDT) on top of frozen TimesFM forecasts to test compute-frugal adaptation.

% \paragraph{E7 (Data-efficiency fine-tuning: full vs PEFT).}
% I compare full fine-tuning to head-only tuning and PEFT methods (LoRA, adapters) across multiple data fractions, and I quantify catastrophic forgetting on a held-out benchmark subset. I explicitly handle the architectural constraint that TimesFM outputs 128 prediction heads by computing loss on the first $H \le 128$ heads when training shorter horizons \cite{timesfm_discussion_214,hu2021lora,houlsby2019adapters}.

% \paragraph{E8 (Embedding transfer: UCR classification).}
% I extract frozen TimesFM embeddings and train lightweight classifiers on UCR datasets, comparing to feature-engineering baselines such as tsfresh \cite{dau2019ucr,christ2018tsfresh}.

% \paragraph{E9 (Anomaly detection).}
% I implement forecast-residual anomaly detection with optional uncertainty-aware scoring and evaluate on NAB and a public anomaly dataset release \cite{nab_lavin2015}.

% \paragraph{E10 (Limited continued pretraining).}
% I perform compute-bounded continued pretraining (full or LoRA-based) on a leakage-free public pretraining set (GiftEvalPretrain) and evaluate gains and forgetting as a function of training steps \cite{gifteval_aksu2024,hu2021lora}.

% \subsection{Comparison table and schedules}
% \begin{table}[!h]
% \centering
% \caption{Summary of experiments (E1--E10).}
% \label{ტაბ:exp_summary}
% \begin{tabular}{p{0.06\linewidth} p{0.34\linewidth} p{0.22\linewidth} p{0.32\linewidth}}
% \hline
% \textbf{Label} & \textbf{Goal} & \textbf{Compute focus} & \textbf{Deliverables} \\
% \hline
% E1 & Zero-shot replication & inference-heavy & tables + DM heatmaps \\
% E2 & Context sweep & inference-heavy & context--MAE curves \\
% E3 & Robustness & inference-heavy & degradation plots \\
% E4 & Calibration & mixed & reliability + CRPS \\
% E5 & XReg covariates & light training & mode comparison \\
% E6 & Residual adapters & light training & gains + diagnostics \\
% E7 & FT vs PEFT + forgetting & training-heavy & data-efficiency + compute frontier \\
% E8 & UCR embeddings & inference-heavy & classification table \\
% E9 & Anomaly detection & rolling inference & NAB scores + PR curves \\
% E10 & Continued pretraining & training-heavy & gain-vs-steps + forgetting \\
% \hline
% \end{tabular}
% \end{table}

% \paragraph{Mermaid Gantt schedules (planning artifacts).}
% \begin{verbatim}
% ```mermaid
% gantt
%   title Low compute plan (single GPU/CPU)
%   dateFormat  YYYY-MM-DD
%   axisFormat  %b %d
%   section Setup
%   Env + logging + manifests      :a1, 2026-03-30, 7d
%   Data + protocol tests          :a2, after a1, 10d
%   section Baselines
%   E1                             :b1, after a2, 10d
%   E2                             :b2, after b1, 7d
%   E3                             :b3, after b2, 14d
%   section Improvements
%   E4                             :c1, after b3, 21d
%   E5                             :c2, after c1, 10d
%   E6                             :c3, after c2, 7d
%   section Training-heavy
%   E7                             :d1, after c3, 35d
%   section Downstream
%   E8                             :e1, after d1, 12d
%   E9                             :e2, after e1, 21d
%   section Continued pretraining
%   E10                            :f1, after e2, 28d
%   section Write-up
%   Writing                         :g1, after f1, 14d
% \end{verbatim}

% \subsection{Reproducibility checklist}
% \begin{itemize}
%     \item I fix random seeds (Python/NumPy/PyTorch) and log them in every run directory.
%     \item I log dataset manifests (dataset name, version/hash, split indices, horizon/context, stride, preprocessing rules).
%     \item I log full TimesFM configuration (checkpoint ID and ForecastConfig fields).
%     \item I store raw predictions and per-series/per-origin losses to enable exact recomputation of all metrics and DM tests.
%     \item I log compute metadata (wall time, GPU-hours, peak VRAM/RAM, batch sizes).
%     \item I track results with W&B and/or TensorBoard and store artifacts (predictions, configs, checkpoints).
% \end{itemize}